\section{Introduction}
\label{sec:introduction}

Large language models (LLMs) are increasingly deployed as autonomous software-engineering agents capable of navigating repositories, interpreting issue descriptions, generating candidate patches, and invoking review loops. Benchmarks such as SWE-bench have accelerated this line of work by grounding evaluation in real GitHub issues and pull requests, where successful repair frequently requires coordinating changes across functions, classes, and files~\cite{jimenez2024swebench}.

This progress also sharpens the cost side of the problem. Recent agent-evaluation work argues that software agents should be evaluated not only by final task success, but also by the tokens, wall-clock time, and infrastructure expense consumed along the trajectory~\cite{fan2025sweeffi}. This paper therefore treats context orchestration as a systems question: can structured retrieval and review improve repository-level software-engineering utility without relying only on larger prompts or more expensive agent loops?

Despite this progress, agentic software-repair systems face a fundamental architectural bottleneck: the trade-off between \textit{context coverage} and \textit{context precision}. When tasked with resolving a complex regression inside a large codebase, an agent needs enough architectural topology to avoid downstream regressions. Long-context models make it tempting to pack large directories, full files, or broad retrieval bundles into a single inference prompt. This can preserve dependencies, but it also creates a secondary failure mode that we call \textit{context dilution}: relevant constraints are buried among unrelated tokens, latency and cost increase, and generated patches become harder to predict.

Conversely, semantic retrieval-augmented generation (RAG) pipelines compress the problem space by extracting isolated chunks or functions using lexical or dense similarity. This preserves a lean token envelope, but it introduces a complementary failure mode that we call \textit{information amputation}. Chunk-based retrieval treats source code as flat text rather than as a dependency-bearing program graph. As a result, it can separate a target function from imported helpers, parent class schemas, configuration decorators, or nearby call contracts that determine whether a patch is valid outside the retrieved snippet.

To study this dichotomy, we introduce the \textit{Context Sphere}, a topology-aware framework that represents repository context as a bounded volume centered at the structural locus of an issue. The framework separates context discovery into two layers: a neural Locator that identifies a candidate centroid, and a deterministic Abstract Syntax Tree (AST) Assembler that expands outward to collect a one-hop neighborhood of local module dependencies. This design aims to retain the high precision of retrieval while recovering the immediate program topology needed for cross-file repair attempts.

The Context Sphere also routes the resulting context volume through an adversarial multi-persona state machine. A Product Manager persona extracts acceptance constraints into a persistent ledger, a Worker persona generates a candidate patch, and a Reviewer persona evaluates the patch against the same bounded neighborhood. When the Reviewer rejects a patch, its diagnostics are appended to the ledger and the Worker receives another attempt. This explicit rejection loop is intended to reduce agentic sycophancy, where a reviewer persona validates an incomplete patch because it inherits the author's assumptions or lacks an independent contract boundary.

We evaluate the current implementation with a 100-case retrieval ablation over pure-Python instances sampled from SWE-bench Verified. Context Sphere achieved 41 successful local repair passes under our \texttt{PASS\_TO\_PASS} smoke/regression harness, compared with 36 passes for a dense Vector baseline and 39 passes for a naive Hybrid retriever. The result shows a 13.8\% relative improvement over dense retrieval at an estimated inference cost of roughly \$0.08 per issue.

Although the empirical stress test in this paper is repository-level bug repair, the architecture is not restricted to repair as an end task. We use repair because it provides an execution-checkable benchmark with concrete pass/fail signals; the same topology-aware retrieval and persona-conditioned orchestration pattern can be applied to broader autonomous software-engineering workflows such as refactoring, dependency analysis, code review, and implementation planning.

This paper makes the following contributions:
\begin{itemize}
    \item We formalize the \textit{Context Sphere}, a hybrid neural-deterministic retrieval strategy that mitigates context dilution and information amputation through AST-guided local topology expansion.
    \item We present a stateful Product Manager--Worker--Reviewer orchestration protocol that converts reviewer rejection into persistent constraint accumulation.
    \item We report a 100-case retrieval ablation comparing dense Vector, Context Sphere, and Hybrid retrieval under a fixed local verification harness.
    \item We observe a retrieval-routing failure mode, the \textit{Hybrid Paradox}, where dense-first context merging can displace structurally critical dependency evidence.
\end{itemize}

While the Context Sphere provides a robust structural foundation for repository-level retrieval, the ablation study identifies a 51-pass theoretical upper bound across retrieval methods, suggesting that further gains require optimizing persona-specific consumption of retrieved context. To address this, we introduce the \textit{Context Projection Model}: a discriminative routing layer that dynamically slices the master Context Sphere into token-efficient, persona-specific \textit{Spheres of Influence}. This extension addresses the context starvation and information bloat observed in baseline multi-agent systems, providing an initial empirical pathway toward raising the current resolution ceiling.
